\section*{Tézy a ciele \typpracerod{} práce}

\normalsize

\subsection*{Tézy}

Tézy sú stručné a jasné východiskové tvrdenia ktoré definujú čo práca rieši, aký je problém, aký je kontext, aké sú hlavné myšlienky alebo prístupy. Tézy teda predstavujú obsahové východiská

Tézy sú teda východiskové myšlienky, ktoré opisujú kontext, problém, prístup alebo teoretické tvrdenia práce.
Nie sú merateľné, sú skôr „čo tvrdím / z čoho vychádzam“.

Príklad téz:

\begin{itemize}
	\item Zabezpečenie komunikácie v počítačových sieťach je kľúčovým prvkom pri ochrane citlivých údajov.
	\item Moderné autentifikačné protokoly umožňujú centralizované riadenie prístupových práv používateľov.
	\item Implementácia bezpečnostných mechanizmov v podnikových sieťach znižuje riziko neoprávneného prístupu.
	\item Analýza reálnych sieťových scenárov poskytuje praktický pohľad na efektivitu jednotlivých bezpečnostných riešení.
\end{itemize}


\subsection*{Ciele}

Ciele sú konkrétne merateľné výstupy, ktoré má práca dosiahnuť, popisujú čo chceme dokázať, čo chceme vytvoriť, čo chceme analyzovať, aké výsledky má práca priniesť. Ciele teda predstavujú konkrétne úlohy vykonané v práci.

Ciele sú teda konkrétne úlohy, ktoré musíš splniť. Sú merateľné, overiteľné a presne definujú, čo bude výsledkom práce.

Príklad cieľov: 

\begin{itemize}
	\item Zanalyzovať dostupné autentifikačné protokoly používané v podnikových sieťach.
	\item Navrhnúť a implementovať model autentifikácie klientov pomocou RADIUS servera.
	\item Otestovať navrhnuté riešenie v simulovanom aj reálnom sieťovom prostredí.
	\item Vyhodnotiť bezpečnostné prínosy a obmedzenia implementovaného riešenia.
\end{itemize}

\newpage